\section{第三阶段 C++20 模型：从 VMA 到 PTE 的可测内存管理器}

虚拟内存不能只靠文字理解。读者至少要实现一个小模型，能把 VMA、PTE、物理页、COW 和缺页处理连起来。这个模型不追求模拟所有硬件细节，而是追求状态可检查：非法地址不能访问，权限收缩必须生效，COW 写入必须复制，引用计数必须闭合。

\subsection{模型边界}

教学模型保留四类对象。

\begin{longtable}{p{0.2\linewidth}p{0.38\linewidth}p{0.3\linewidth}}
\toprule
对象 & 责任 & 不做什么 \\
\midrule
\code{PhysicalMemory} & 分配页框、维护引用计数、复制页 & 不模拟 NUMA、cache、DMA \\
\code{AddressSpace} & 管理 VMA 和页表 & 不模拟多级页表的每级页表页 \\
\code{PageTableEntry} & 保存页框、权限、COW 标志 & 不模拟 accessed/dirty 的硬件置位细节 \\
\code{FaultHandler} & 按 VMA/PTE/访问类型处理缺页 & 不模拟真实信号栈和架构异常帧 \\
\bottomrule
\end{longtable}

先省略一部分硬件细节，是为了让核心不变量清晰。等这些不变量稳定后，再加入 accessed/dirty、swap、file mapping、TLB shootdown 和 memcg，才不会把系统写成一团不可测代码。

\subsection{类型和常量}

先用强类型别名和命名常量写清地址单位。所有有语义的数字都要命名。

\begin{lstlisting}[language=C++]
#include <algorithm>
#include <cstdint>
#include <deque>
#include <limits>
#include <stdexcept>
#include <unordered_map>
#include <utility>
#include <vector>

using VirtualAddress = std::uint64_t;
using PageNumber = std::uint64_t;
using FrameNumber = std::uint64_t;

constexpr std::uint64_t VM_PAGE_SIZE = 4096;
constexpr VirtualAddress VM_USER_TOP = 0x0000800000000000ULL;
constexpr FrameNumber VM_INVALID_FRAME =
    std::numeric_limits<FrameNumber>::max();

enum class AccessKind {
  Read,
  Write,
  Execute,
};
\end{lstlisting}

地址计算要集中放在小函数里，避免到处手写除法和取模。

\begin{lstlisting}[language=C++]
constexpr PageNumber page_number(VirtualAddress address) {
  return address / VM_PAGE_SIZE;
}

constexpr std::uint64_t page_offset(VirtualAddress address) {
  return address % VM_PAGE_SIZE;
}

constexpr VirtualAddress page_down(VirtualAddress address) {
  return address - page_offset(address);
}

constexpr VirtualAddress page_up(VirtualAddress address) {
  const auto offset = page_offset(address);
  if (offset == 0) {
    return address;
  }
  return address + (VM_PAGE_SIZE - offset);
}
\end{lstlisting}

\subsection{权限和 VMA}

VMA 是区间策略，不是当前硬件映射。它说明地址是否属于进程、理论上允许什么访问、缺页时应该怎样补页。

\begin{lstlisting}[language=C++]
struct VmPermissions {
  bool readable = false;
  bool writable = false;
  bool executable = false;
};

bool allows(const VmPermissions& permissions, AccessKind access) {
  switch (access) {
    case AccessKind::Read:
      return permissions.readable;
    case AccessKind::Write:
      return permissions.writable;
    case AccessKind::Execute:
      return permissions.executable;
  }
  throw std::logic_error("unknown access kind");
}

enum class VmaKind {
  Anonymous,
  FilePrivate,
  FileShared,
};

struct Vma {
  VirtualAddress start = 0;
  VirtualAddress end = 0;
  VmPermissions permissions;
  VmaKind kind = VmaKind::Anonymous;
  std::uint64_t file_offset = 0;
};
\end{lstlisting}

VMA 插入时必须检查页对齐、非空区间、用户地址范围和不重叠。教学模型可用有序 vector；查找是线性扫描，语义最直观。

\begin{lstlisting}[language=C++]
class VmaTable {
 public:
  void insert(Vma vma) {
    this->validate_new_vma(vma);
    auto position = std::ranges::lower_bound(
        this->vmas_, vma.start, {}, &Vma::start);
    this->vmas_.insert(position, std::move(vma));
  }

  [[nodiscard]] const Vma* find(VirtualAddress address) const {
    for (const auto& vma : this->vmas_) {
      if (address < vma.start) {
        return nullptr;
      }
      if (vma.start <= address && address < vma.end) {
        return &vma;
      }
    }
    return nullptr;
  }

 private:
  void validate_new_vma(const Vma& vma) const {
    if (vma.start >= vma.end) {
      throw std::invalid_argument("empty VMA");
    }
    if (page_offset(vma.start) != 0 || page_offset(vma.end) != 0) {
      throw std::invalid_argument("VMA must be page aligned");
    }
    if (vma.end > VM_USER_TOP) {
      throw std::invalid_argument("VMA outside user range");
    }
    for (const auto& existing : this->vmas_) {
      const bool disjoint = vma.end <= existing.start || existing.end <= vma.start;
      if (!disjoint) {
        throw std::invalid_argument("overlapping VMA");
      }
    }
  }

  std::vector<Vma> vmas_;
};
\end{lstlisting}

这段代码故意没有实现自动合并。教学阶段先让每次插入的边界完全显式；等测试覆盖充分后，再把相邻 VMA 合并作为优化加入。

\subsection{物理页和引用计数}

物理内存模型用 frame number 标识页框，页内容用固定大小数组表示。为了节省书中篇幅，可以把内容简化为 byte vector；真实实现中应避免在热路径频繁分配。

\begin{lstlisting}[language=C++]
struct PhysicalPage {
  std::vector<std::uint8_t> bytes = std::vector<std::uint8_t>(VM_PAGE_SIZE);
  std::uint32_t refcount = 0;
};

class PhysicalMemory {
 public:
  FrameNumber allocate_zeroed() {
    FrameNumber frame = VM_INVALID_FRAME;
    if (!this->free_list_.empty()) {
      frame = this->free_list_.front();
      this->free_list_.pop_front();
      auto& page = this->pages_.at(frame);
      std::ranges::fill(page.bytes, 0);
      page.refcount = 1;
      return frame;
    }

    frame = static_cast<FrameNumber>(this->pages_.size());
    PhysicalPage page;
    page.refcount = 1;
    this->pages_.push_back(std::move(page));
    return frame;
  }

  FrameNumber clone(FrameNumber source) {
    const FrameNumber target = this->allocate_zeroed();
    this->pages_.at(target).bytes = this->pages_.at(source).bytes;
    return target;
  }

  void retain(FrameNumber frame) {
    this->pages_.at(frame).refcount += 1;
  }

  void release(FrameNumber frame) {
    auto& page = this->pages_.at(frame);
    if (page.refcount == 0) {
      throw std::logic_error("double release of physical page");
    }
    page.refcount -= 1;
    if (page.refcount == 0) {
      this->free_list_.push_back(frame);
    }
  }

  [[nodiscard]] std::uint32_t refcount(FrameNumber frame) const {
    return this->pages_.at(frame).refcount;
  }

 private:
  std::vector<PhysicalPage> pages_;
  std::deque<FrameNumber> free_list_;
};
\end{lstlisting}

引用计数是 COW 的核心证据。每个测试结束时，物理页活跃数、引用计数和页表映射数要能对上，否则模型可能泄漏或重复释放。

\subsection{PTE 和地址空间}

PTE 保存当前硬件可见映射。VMA 允许写不等于 PTE 当前可写；COW 页的 VMA 允许写，但 PTE 要暂时只读，让第一次写触发 fault。

\begin{lstlisting}[language=C++]
struct PageTableEntry {
  FrameNumber frame = VM_INVALID_FRAME;
  VmPermissions permissions;
  bool present = false;
  bool cow = false;
};

class AddressSpace {
 public:
  void map_vma(Vma vma) {
    this->vmas_.insert(vma);
  }

  [[nodiscard]] const Vma* find_vma(VirtualAddress address) const {
    return this->vmas_.find(address);
  }

  [[nodiscard]] PageTableEntry* find_pte(VirtualAddress address) {
    const auto iterator = this->page_table_.find(page_number(address));
    if (iterator == this->page_table_.end()) {
      return nullptr;
    }
    return &iterator->second;
  }

  void install_pte(VirtualAddress address, PageTableEntry pte) {
    this->page_table_[page_number(address)] = pte;
  }

 private:
  VmaTable vmas_;
  std::unordered_map<PageNumber, PageTableEntry> page_table_;
};
\end{lstlisting}

这个页表是哈希表，不是多级页表。它足以表达“某个虚拟页当前是否 present、权限是什么、是否 COW”。后续实践卷再把它替换成四级页表或模拟 TLB。

\subsection{缺页处理}

缺页处理先查 VMA，再查权限，再决定是新分配、COW 还是权限错误。返回结果要区分用户可见错误。

\begin{lstlisting}[language=C++]
enum class FaultResult {
  Resolved,
  SegmentationFault,
  ProtectionFault,
};

class FaultHandler {
 public:
  explicit FaultHandler(PhysicalMemory& memory) : memory_(memory) {}

  FaultResult handle(AddressSpace& space, VirtualAddress address, AccessKind access) {
    const Vma* vma = space.find_vma(address);
    if (vma == nullptr) {
      return FaultResult::SegmentationFault;
    }
    if (!allows(vma->permissions, access)) {
      return FaultResult::ProtectionFault;
    }

    PageTableEntry* pte = space.find_pte(address);
    if (pte != nullptr && pte->present) {
      return this->handle_present_pte(*pte, access);
    }

    PageTableEntry new_pte;
    new_pte.frame = this->memory_.allocate_zeroed();
    new_pte.present = true;
    new_pte.permissions = vma->permissions;
    space.install_pte(address, new_pte);
    return FaultResult::Resolved;
  }

 private:
  FaultResult handle_present_pte(PageTableEntry& pte, AccessKind access) {
    if (allows(pte.permissions, access)) {
      return FaultResult::Resolved;
    }
    if (access == AccessKind::Write && pte.cow) {
      return this->resolve_cow(pte);
    }
    return FaultResult::ProtectionFault;
  }

  FaultResult resolve_cow(PageTableEntry& pte) {
    if (this->memory_.refcount(pte.frame) == 1) {
      pte.permissions.writable = true;
      pte.cow = false;
      return FaultResult::Resolved;
    }

    const FrameNumber old_frame = pte.frame;
    const FrameNumber new_frame = this->memory_.clone(old_frame);
    this->memory_.release(old_frame);
    pte.frame = new_frame;
    pte.permissions.writable = true;
    pte.cow = false;
    return FaultResult::Resolved;
  }

  PhysicalMemory& memory_;
};
\end{lstlisting}

这段实现有一个清晰的不变量：COW 处理前，PTE 不可写；处理后，当前地址空间获得可写私有页；旧 frame 引用计数减少。若 refcount 为 1，不复制，只把权限改回可写。

\subsection{fork 的 COW 准备}

fork 时复制页表，但共享物理页，把双方 PTE 改为只读 COW。VMA 不需要复制物理页，只要复制区间策略。

\begin{lstlisting}
fork address space:
  copy VMA table
  for each present writable anonymous PTE:
    clear writable in parent PTE
    set cow in parent PTE
    create child PTE pointing to same frame
    clear writable in child PTE
    set cow in child PTE
    retain physical frame
\end{lstlisting}

教学代码实现 fork 时要特别注意父进程 PTE 也要改成只读。如果只改子进程，父进程写入不会 fault，会直接污染共享页。真实系统还要 flush 父进程对应 TLB，否则父 CPU 可能继续使用旧 writable 权限。

\subsection{mprotect 的模型}

\code{mprotect} 至少要改 VMA 权限和已有 PTE 权限。教学模型可以先拒绝跨 VMA 的复杂修改，只允许整段 VMA 修改；实践卷再实现切分。

\begin{lstlisting}
mprotect simplified:
  find exact VMA [start, end)
  verify new permissions do not exceed may permissions
  update VMA permissions
  for each PTE in range:
    pte.permissions = intersect(pte.permissions, new_permissions)
    if new permissions remove write:
      record need_tlb_flush
\end{lstlisting}

注意“扩大权限”不能只改 PTE，还要检查原始 \code{may} 权限；“缩小权限”不能只改 VMA，还要改已经 present 的 PTE。否则已有映射会继续按旧权限访问。

\subsection{测试用例}

这个模型至少要有以下测试。

\begin{longtable}{p{0.24\linewidth}p{0.36\linewidth}p{0.28\linewidth}}
\toprule
测试 & 步骤 & 期望 \\
\midrule
未映射地址 & 不建 VMA，访问地址 & 返回段错误结果 \\
懒分配 & 建匿名 VMA，首次读 & 分配零页并 resolved \\
写只读 VMA & VMA 只读，执行写 fault & 返回权限错误结果 \\
COW refcount>1 & 父子共享页，子写 & 子获得新 frame，旧 frame refcount 减一 \\
COW refcount=1 & 单独持有 COW 页后写 & 不复制，只恢复 writable \\
mprotect 收缩 & RW 改 R 后写 & 写失败，旧权限不能生效 \\
double release & 对同一 frame release 两次 & 抛出逻辑错误或测试失败 \\
\bottomrule
\end{longtable}

测试要断言内部状态，而不是只断言返回值。例如 COW 测试必须检查父子 frame 是否不同、内容是否正确、旧 frame 引用计数是否符合预期。只有这样才能发现“看似成功但泄漏引用”的错误。

\subsection{从教学模型到真实内核}

这个模型和真实内核差距很大，但差距是有方向的。

\begin{itemize}
\tightlist
\item 把哈希页表替换为多级页表，可学习页表页分配和页表遍历。
\item 给 PTE 加 accessed/dirty，可学习回收和写回。
\item 给 VMA 加文件对象和 offset，可学习 mmap 与 page cache。
\item 给地址空间加 active CPU 集合，可学习 TLB shootdown。
\item 给物理页加 memcg、zone、order 和 migratetype，可学习 buddy、NUMA、compaction。
\item 给 fault handler 加可睡眠点和 retry，可学习真实缺页路径的并发。
\end{itemize}

不要一开始就实现所有功能。先让最小模型通过严格测试，再逐步加入复杂字段。操作系统工程的质量来自可维护的不变量，而不是一次性写出大量不可验证的机制。
